% ===========================================================================
%  user/07-community/00-chapter.tex
% ===========================================================================
\chapter{Community}\label{cap:community}

It is the area where users meet, and the only one in the application that
needs an account. In the menu it sits under \voce{Share}; the page has four
tabs:

\begin{center}
\begin{tabular}{@{}L{34mm}L{104mm}@{}}
\toprule
\textbf{Tab} & \textbf{What it is for} \\
\midrule
\rigaTenue
\voce{Marketplace} & Buying and selling strategies, approaches, overlays and
  weighting methods. It is the subject of this chapter. \\
\voce{Messages} & Conversations with other users, for example with the
  author of a strategy. \\
\rigaTenue
\voce{Feedback} & The reports and suggestions sent to the Lyra team. \\
\voce{My account} & Profile, security and subscription
  (\vediCap{account}). \\
\bottomrule
\end{tabular}
\end{center}

\figuraSchermata{1.0}{figure/schermate/marketplace.png}
  {The \voce{Community} without an account. The rest of the application keeps
   working: the page says so.}
  {community-signed-out}

\section{Why an account is needed here and not elsewhere}\label{sez:why-account}

Because here there is someone on the other side. Building and testing
strategies happens entirely on your own computer, and asking for an account
to do it would make no sense. Buying, selling and granting a licence require
the system to know who is who.

\section{Buying}\label{sez:buy}

\begin{pratica}{From the storefront to the system}
\begin{enumerate}
  \item Open the artifact's page: description, metrics declared by the
        seller, comments, sometimes a video.
  \item \pulsante{Buy} --- if it is free, access is immediate; if it is paid,
        the payment flow opens.
  \item Once the payment is done you go back to the application. If access
        does not show as active yet, \textbf{it is not an error}: the
        confirmation arrives a moment later. The page says so and asks you to
        wait.
  \item \pulsante{Install} --- the artifact is written to the system, and
        from then on it appears in the lists like all the others.
\end{enumerate}
\end{pratica}

\begin{nota}{Before installing, the conflicts}
If something with the same name already exists in the system, the
application says so \emph{before} writing. It is worth reading: an
installation that overwrites a strategy you were working on is lost work.
\end{nota}

\subsection{If a purchase does not show up}

There is a command for exactly this, on the artifact's page, that explains
why it does not show as purchased. It tells apart two things that at first
sight look the same: \enquote{you did not buy it} and \enquote{I could not
ask}. They are opposites, and the second is solved by trying again.

If the payment went through but access does not open, there is a
reconciliation command: it checks the payment again with the provider and,
only if it shows as collected, completes the operation.

\section{Selling}\label{sez:sell}

To publish you need two things, in this order:

\begin{enumerate}
  \item \textbf{The marketplace must be open.} During the beta it is closed:
        nobody uploads anything. Items already in the storefront stay
        visible.
  \item \textbf{A paid plan.} With the marketplace open, any paid plan will
        do, from Pro up (\vediCap{account}). On Free you can browse, download
        and buy, but not sell.
\end{enumerate}

\begin{perche}{Why a subscription even to publish for free}
It is not about the money. Whoever publishes puts their name next to a
strategy that others will run with their own money: a filter at the entrance
is worth more than any check after the fact. Someone without a subscription
remains a full user --- they browse, download, buy and use. They just cannot
put things in the storefront.
\end{perche}

\noindent
A published artifact goes online at once: there is no approval queue. The
description and the video can be corrected later; the bundle of files
cannot.

\section{If an artifact does not perform as declared}\label{sez:dispute}

You open a dispute from the artifact's page. It is not a discussion: the
system runs two automatic checks.

\begin{center}
\begin{tabular}{@{}L{44mm}L{94mm}@{}}
\toprule
\textbf{Check} & \textbf{What it compares} \\
\midrule
\rigaTenue
\voce{Integrity} & The installed files with those sold, one by one. \\
\voce{Performance} & The metrics declared by the seller with those obtained
  by running the calculation again on your own system. \\
\bottomrule
\end{tabular}
\end{center}

\noindent
The case is a conversation between buyer, seller and the administration,
with the results of the two checks attached. It can be withdrawn at any
time.

\begin{nota}{Two different results do not mean a deception}
A strategy can perform differently because the data are different: a
different universe of instruments, a different period, different costs. The
check is there to establish \emph{by how much} they differ, which is the
starting point of any sensible discussion.
\end{nota}
